%% kapa: exam=1 round=30 year=2019 period=2 form=A subject=appraisal_law subject_ko=감정평가관계법규 seq=14 number=14
%% source: Q-Net 감정평가사 기출문제 (cst00309) artlSeq=5205205 file=extracted/1008/5205205/2019년 제30회 감정평가사 1차시험 문제지/2019년 제30회 감정평가사 1차시험 문제지/2019년도 제30회 감정평가사 1차 2교시 A형.hwp
%% conversion: kordoc 4.24.51 -> tools/gen1.py (auto); review=proofread-sonnet; flags=none
\begin{question}{14}
감정평가 및 감정평가사에 관한 법령상 감정평가법인에 관한 설명으로 옳은 것은?
\begin{choices}
  \item 감정평가법인의 주사무소 및 분사무소에 주재하는 감정평가사가 각각 3명이면 설립 기준을 충족하지 못한다.
  \item 감정평가법인을 해산하려는 경우에는 국토교통부장관의 인가를 받아야 한다.
  \item 감정평가법인은 사원 전원의 동의 또는 주주총회의 의결이 있는 때에는 국토교통부장관의 인가를 받아 다른 감정평가법인과 합병할 수 있다.
  \item 자본금 미달은 감정평가법인의 해산 사유에 해당한다.
  \item 국토교통부장관은 감정평가업자가 업무정지처분 기간 중에 법원에 계속 중인 소송 또는 경매를 위한 토지등의 감정평가업무를 한 경우 가중하여 업무의 정지를 명할 수 있다.
\end{choices}
\end{question}
\answer{3}
